\section*{Chapitre \ref{ch:b3:colloids} --- Interfaces, tensioactifs et colloïdes}

\begin{solution}{exo:b3:colloids:1}
$r = \qty{0.50}{\micro m}$~: $\Delta p = 2 \times 0.0720/\num{5.0e-7} = \qty{2.9e5}{Pa}$~; à l'intérieur, $1.0 + 2.9 = \qty{3.9}{bar}$.
\end{solution}

\begin{solution}{exo:b3:colloids:2}
$\ln(p/p^*) = 2 \times 0.0720 \times \num{1.807e-5}/(\num{1.0e-8} \times 8.314 \times 298.15) = 0.105$~; $p/p^* = 1.11$.
\end{solution}

\begin{solution}{exo:b3:colloids:3}
$\cos\theta = (40 - 20)/72 = 0.278$, $\theta = \ang{74}$~: au-dessous de
$90^\circ$, le liquide \omterm{def:b3:colloids:wetting}{mouille} partiellement le solide.
\end{solution}

\begin{solution}{exo:b3:colloids:4}
NaCl~: $I = \qty{10}{mol.m^{-3}}$, $\kappa^{-1} = 0.96\sqrt{100/10} = \qty{3.0}{nm}$. \ce{MgSO4}~: $I = \qty{40}{mol.m^{-3}}$, \qty{1.5}{nm}.
\end{solution}

\begin{solution}{exo:b3:colloids:5}
$\Gamma = \num{12.6e-3}/(2 \times 8.314 \times 298.15) = \qty{2.54e-6}{mol.m^{-2}}$~; aire $1/(\Gamma N_A) =
\qty{0.65}{nm^2}$ par molécule.
\end{solution}

\begin{solution}{exo:b3:colloids:6}
Au-dessous de la CMC, la pente vaut $-12.9$ \unit{mN/m} par unité de
$\log c$ (de $-5.0$ à $-4.0$)~; le
palier se situe à 33.4. La droite l'atteint en $\log c = -4.0 + (39.1 - 33.4)/12.9 = -3.56$~:
CMC
$\approx \qty{2.8e-4}{mol/L}$.
\end{solution}

\begin{solution}{exo:b3:colloids:7}
SDS~: $P = 0.350/(0.60 \times 1.67) = 0.35$, à la frontière entre sphères
et cylindres courts~: \omterm{def:b3:colloids:micelle}{micelles} sphériques au voisinage de la CMC.
Lipide~: $P = 0.700/(0.65 \times 1.67) = 0.64$~:
bicouches, donc vésicules.
\end{solution}

\begin{solution}{exo:b3:colloids:8}
\ce{CaCl2}~: $60/64 = \qty{0.94}{mM}$ de \ce{Ca^{2+}}~; \ce{AlCl3}~: $60/729 = \qty{0.082}{mM}$ de \ce{Al^{3+}}. \ce{Na3PO4}
coagule par son \ce{Na+} (vers \qty{20}{mM} de sel)~; le phosphate, un
co-ion, peut même s'adsorber et augmenter la charge négative.
\end{solution}

\begin{solution}{exo:b3:colloids:9}
Partie hydrophile~: $6 \times 44.0 + 17.0 = \qty{281.0}{g/mol}$ sur
\qty{450.0}{g/mol} (masses atomiques du livre)~:
HLB $= 20 \times 281.0/450.0 = 12.5$, un émulsifiant huile dans eau.
\end{solution}

\begin{solution}{exo:b3:colloids:10}
L'huile des petites gouttelettes est plus soluble dans la phase continue
(d'un facteur $\exp(2\gamma V_{\mathrm m}/rRT)$, dont l'exposant est dix
fois plus grand pour \qty{50}{nm} que pour \qty{500}{nm})~: elle
diffuse vers les grosses gouttelettes, qui grossissent tandis que les
petites disparaissent. Ce \omterm{def:b3:colloids:ripening}{mûrissement d'Ostwald} fait grossir une \omterm{def:b3:colloids:colloid}{émulsion}
même si deux gouttelettes ne fusionnent jamais.
\end{solution}

\begin{solution}{exo:b3:colloids:11}
Le ménisque est une calotte sphérique de rayon $r/\cos\theta$~: le liquide
juste au-dessous est à une pression plus basse de $2\gamma\cos\theta/r$,
et la colonne monte jusqu'à ce que $\rho gh$ compense. Pour l'eau,
$h = 2 \times 0.072/(997 \times 9.81 \times \num{1.0e-4}) = \qty{0.15}{m}$.
\end{solution}

\begin{solution}{exo:b3:colloids:12}
À \qty{100}{mM}, la \omterm{def:b3:electrode-kinetics:double-layer}{longueur de Debye} est inférieure à \qty{1}{nm} et la
répulsion s'éteint avant l'attraction~: il ne reste aucune barrière
(au-delà de la courbe à \qty{50}{mM}). Un polymère libre n'agit pas à la
surface~; une couche de polymère greffé repousse au contact quel que soit
le sel, par l'entropie perdue quand les chaînes se recouvrent.
\end{solution}

\begin{solution}{pb:b3:colloids:1}
\textbf{1.} Des substitutions dans le réseau cristallin (\ce{Al^{3+}} à
la place de \ce{Si^{4+}}, \ce{Mg^{2+}} à la place de \ce{Al^{3+}})
laissent une charge négative permanente, compensée par des cations
échangeables.
\textbf{2.} $I = \frac12(1.0 + 1.0) = \qty{1.0}{mM}$.
\textbf{3.} $\kappa^{-1} = \qty{9.6}{nm}$.
\textbf{4.} $a\kappa = 10$~: la double couche est mince par rapport à la
particule, condition de l'approximation de Derjaguin.
\textbf{5.} $v = 2 \times 1650 \times 9.81 \times (\num{1.0e-7})^2/(9 \times \num{0.89e-3}) = \qty{4.0e-8}{m/s}$,
environ \qty{3.5}{mm} par jour.
\textbf{6.} Leurs \omterm{def:b3:electrode-kinetics:double-layer}{couches diffuses} se repoussent avant que l'attraction
de van der Waals l'emporte.
\textbf{7.} $V = 2\pi\varepsilon_r\varepsilon_0a\psi^2\ln(1 + \eu^{-\kappa h}) - Aa/12h$.
\textbf{8.} Environ $39\,kT$.
\textbf{9.} Une fraction de l'ordre de $\eu^{-39} \approx 10^{-17}$ des
collisions~: pratiquement jamais.
\textbf{10.} Elle disparaît~: toute collision conduit au contact.
\textbf{11.} Une \omterm{def:b3:electrode-kinetics:double-layer}{longueur de Debye} plus courte confine la répulsion aux
petites distances, où l'attraction, qui croît comme $1/h$, domine.
\textbf{12.} La CCC varie comme $z^{-6}$ en fonction de la charge du
contre-ion.
\textbf{13.} $50/64 = \qty{0.78}{mM}$.
\textbf{14.} $50/729 = \qty{0.069}{mM}$.
\textbf{15.} Les cations sont les contre-ions des particules négatives~:
ils s'accumulent dans la double couche et l'écrantent.
\textbf{16.} Le sulfate est un co-ion, repoussé par les particules.
\textbf{17.} \qty{0.069}{mol} de \ce{Al^{3+}} par mètre cube.
\textbf{18.} \qty{0.034}{mol} de $\ce{Al2(SO4)3}$.
\textbf{19.} $0.0343 \times 342.3 = \qty{11.7}{g}$.
\textbf{20.} \qty{11.7}{kg} par heure.
\textbf{21.} $\ce{Al^3+ + 3 HCO3- -> Al(OH)3 + 3 CO2}$.
\textbf{22.} $3 \times 0.069 = \qty{0.21}{mM}$ sur \qty{1.0}{mM}~: un
cinquième. L'hydrogénocarbonate restant et le \ce{CO2} dissous tamponnent
l'eau~: le pH ne baisse que modérément.
\textbf{23.} Les espèces de l'aluminium très chargées s'adsorbent et
inversent la charge des particules, qui se repoussent de nouveau.
\textbf{24.} \textbf{Environ \qty{12}{g} de sulfate d'aluminium par mètre
cube} (\qty{11.7}{g}), avant le supplément nécessaire en pratique pour
l'hydrolyse.
\end{solution}
